% ID: SEL-D0031-OB-E01
% ESTADO: validada. S032, cotejada con examen y criterios.
Sean $x,y$ las alfombras de seda y lana. Las restricciones son
$x,y\ge0$, $2x+3y\le600$ para el trabajo manual y $2x+y\le480$
para el trabajo de máquina. El beneficio es $150x+100y$.

Sobre los ejes, los límites más restrictivos son $x=240$ e $y=200$.
Al restar las dos ecuaciones frontera se obtiene $2y=120$,
luego su intersección es $(210,60)$.
\begin{center}
\setlength{\unitlength}{7mm}
\begin{picture}(12,9)(-2,-1)
\qbezier(0,0)(4.5,0)(9,0)\put(9.2,0){$x$}
\qbezier(0,0)(0,3.6)(0,7.2)\put(0,7.3){$y$}
\qbezier(0,6.667)(3.5,4.3335)(7,2)
\qbezier(7,2)(7.5,1)(8,0)
\put(-0.2,6.67){\makebox(0,0)[r]{$(0,200)$}}
\put(7.2,2.2){$(210,60)$}
\put(8,-0.5){\makebox(0,0)[t]{$(240,0)$}}
\end{picture}
\end{center}
El beneficio toma los valores $0,36000,37500,20000$ en
$(0,0)$, $(240,0)$, $(210,60)$ y $(0,200)$. Por linealidad,
el máximo se alcanza en el mejor vértice, que además es entero.
Se deben producir \textbf{210 alfombras de seda y 60 de lana}.
El beneficio máximo es $\boxed{37500}$ euros.
